\documentclass[11pt]{article}
\usepackage[T1]{fontenc}
\usepackage{lmodern,amsmath,amssymb,amsthm,booktabs}
\usepackage[margin=1in]{geometry}
\usepackage[colorlinks=true,linkcolor=blue,urlcolor=blue,citecolor=blue]{hyperref}
\newtheorem{theorem}{Theorem}
\newtheorem{lemma}{Lemma}
\newcommand{\PrE}{\Pr}
\newcommand{\repo}{https://github.com/Aafimalek/novel-discoveries/tree/research/parallel-foundations/width-three-posets}
\title{A balanced pair for a two-by-three branch in a poset}
\author{Aafi Khan Malek\\[-2pt]
\small Independent researcher, Ahmedabad, India\\[-2pt]
\small \href{mailto:aafimalek2032@gmail.com}{aafimalek2032@gmail.com}}
\date{27 November 2026}
\begin{document}
\maketitle

\begin{abstract}
We prove a local sufficient condition for the existence of a balanced pair
in a finite poset. If a specified exclusive upset has a root and ordered
arms of lengths two and three, a balanced pair is guaranteed when
the probability that the distinguished element precedes the root is at most
$7/30$, subject to a predecessor-set hypothesis. With one additional
cross-arm relation, the sufficient threshold is $1/4$. A finite table of
integer inequalities is checked in Lean 4, together with the reduction from
linear extensions and the full restricted theorem. A cutoff-mixture example
shows that $7/30$ is sharp for the counting relaxation, not necessarily for posets.
\end{abstract}

\section{Statement}
For a finite poset $P$, let $e(P)$ be its number of linear extensions.
Write $u\prec_L v$ when $u$ precedes $v$ in an extension $L$, and write
$\PrE(u\prec v)$ for the probability under the uniform choice of $L$.
A pair of distinct elements is balanced if that probability belongs to
$[1/3,2/3]$. The $1/3$--$2/3$ conjecture asks for a balanced pair in every
finite non-chain; for context on the width-two case see \cite{sah}.

\begin{theorem}\label{thm:main}
Let $a,b,x,X,y,Y,Z$ be distinct elements of $P$, and suppose that
$b\not\le_P a$ and every strict predecessor of $a$ is a strict
predecessor of $b$. Suppose further that
\[
S=\{z\in P:b\le_P z,\ a\not<_P z\}
  =\{b,x,X,y,Y,Z\},\qquad
b<_P x<_P X,\quad b<_P y<_P Y<_P Z.
\]
Set $p=\PrE(a\prec b)$. There is a balanced pair among these seven
elements if $p\le7/30$, or if $x<_P Z$ and $p\le1/4$.
Other cross-arm comparisons are allowed.
\end{theorem}

\section{From extensions to cutoff mixtures}
There are ten ways to interleave $(x,X)$ and $(y,Y,Z)$ after $b$.
Use $\sigma$ for such an induced order of $S$ and, for $0\le j\le6$,
let $q_{\sigma,j}$ count extensions with induced order $\sigma$ in
which exactly $j$ members of $S$ precede $a$.

\begin{lemma}\label{lem:monotone}
For every $\sigma$, $q_{\sigma,j+1}\le q_{\sigma,j}$ for $0\le j<6$.
\end{lemma}
\begin{proof}
For an extension counted at $j+1$, exchange $a$ with the last preceding
member $u$ of $S$. Moving $a$ earlier is legal: each of its strict
predecessors precedes $b$ in the poset and thus precedes $u$.
Moving $u$ later is legal as well. If a strict successor $z$ of $u$
occurred between $u$ and $a$, then $b\le_P z$. Moreover $a\not<_P z$,
since $z$ preceded $a$ in the extension. Thus $z\in S$, contrary to
the choice of $u$. The exchange preserves $\sigma$ and decreases the
cut by one. The first member of $S$ after $a$ in the resulting extension
is $u$, so the exchange can be reversed. This gives the injection.
\end{proof}

Write $q_{\sigma,7}=0$ and $w_{\sigma,k}=q_{\sigma,k}-q_{\sigma,k+1}$
for $0\le k\le6$. Thus the histogram is a nonnegative integer mixture
of 70 components. Component $(\sigma,k)$ consists of one occurrence
of each cut $j=0,\dots,k$ and has mass $k+1$.
Let $C_{\sigma,k}(u,v)$ count its cuts with $u\prec v$. For $u,v\in S$
this is either $0$ or $k+1$, according to $\sigma$. For $v\in S$ of
zero-based rank $r_\sigma(v)$,
\[
 C_{\sigma,k}(a,v)=\min\{k+1,r_\sigma(v)+1\},
 \qquad C_{\sigma,k}(v,a)=k+1-C_{\sigma,k}(a,v).
\]
The extension count for a comparison is therefore
$C(u,v)=\sum_{\sigma,k}w_{\sigma,k}C_{\sigma,k}(u,v)$ and the total
mass is $N=\sum_{\sigma,k}w_{\sigma,k}(k+1)=e(P)$.
If $x<_P Z$, all weights on orders with $Z$ before $x$ vanish.

\section{Eliminating strict-majority orders}
Assume no pair among the seven is balanced. Direct each pair in the
direction whose extension fraction is strictly greater than $2/3$.
A directed triangle is impossible: by the union bound, its three
events would have a common extension, whereas no linear order can
realize a directed cycle. A tournament with no directed triangle is
transitive, since $u$ before $v$ and $v$ before $w$ force $u$ before
$w$. The resulting total order begins with $b$: it precedes the five
arm elements in every extension, and $\PrE(b\prec a)=1-p>2/3$
under either hypothesis. There are 60 possible such orders, obtained
from the ten interleavings and the six placements of $a$ after $b$.

For a nonnegative list of comparison coefficients $c=(c_{uv})$,
let $\Phi_c(C)=\sum c_{uv}C(u,v)$, $B=c_{ba}$, and
$T=\sum_{uv\ne ba}c_{uv}$. Table~\ref{tab:cert} lists 25 coefficient
vectors. The exact checks in the supplement show:
\begin{enumerate}
\item For each listed row and each of the 70 components,
\[
 30\Phi_c(C_{\sigma,k})\le (23B+20T)(k+1),\qquad T>0.
\]
\item Every possible majority order directs all the comparisons with
positive coefficient in at least one row forward. If $a$ is not
immediately after $b$, one may choose a row with $B=0$.
\end{enumerate}
For the $7/30$ bound, sum the first inequality over the mixture and
choose a row as in item 2. Its $ba$ term contributes at least
$23BN/30$ and its other terms strictly more than $20TN/30$.
This contradicts the summed inequality.

For the restricted $1/4$ bound, the same rows with $B=0$ dispose
of the 50 later placements. With $a$ immediately after $b$, use
one of the following coefficient vectors, according to which arm
begins first in the majority order:
\[
 c_X=12\,ba+6\,ax+2\,xy+xY,\qquad
 c_Y=12\,ba+6\,ay+2\,yx+yX.
\]
Here a symbol such as $xy$ denotes the comparison $x\prec y$.
Each active component satisfies $\Phi_c\le15(k+1)$; for $c_X$ the
check uses only orders with $x$ before $Z$, and for $c_Y$ it uses
all orders. In either case $C(b,a)\ge 3N/4$ and the other three
comparisons are strong. Hence $\Phi_c>12(3N/4)+9(2N/3)=15N$,
a contradiction. This proves Theorem~\ref{thm:main}.

\section{A sharp obstruction in the counting model}
The relaxation alone cannot raise $7/30$ even to $1/4$ without
the extra relation. In the order $bxXyYZ$ put weight $5s$ on each
cut $0$ through $3$, and additional weight $t$ on cuts $0,1$.
In the order $byYZxX$ put weight $2s$ on each cut $0$ through $4$.
For $0<t<9s$ the majority order is $baxXyYZ$, with every forward
comparison strictly above $2/3$; the total mass is $30s+2t$ and
$p=(7s+t)/(30s+2t)$. With $s=2,t=1$ this gives $15/62$ and
no balanced comparison. The probabilities can approach $7/30$
from above. For each forward pair the excess of three times its count
over twice the mass is linear in $s,t$. The audit checks it is
nonnegative at $t=0$ and $t=9s$, with at least one strict endpoint;
this proves strict positivity throughout the indicated open interval.
These are histograms, not a construction of a poset.

\section{Reproducibility}
The \href{\repo/verification/two-by-three-balanced}{Lean package}
proves \texttt{WidthThree.balanced\_of\_twoByThreeShape}. It checks
the actual extension counts, their antitone histogram, the exact
cutoff decomposition, the strict-majority order classification,
and the finite inequalities. The replay used Lean 4.36.0 and the
pinned Mathlib in the package. The theorem audit lists only
\texttt{propext}, \texttt{Classical.choice}, and \texttt{Quot.sound}.
An independent short Python audit accompanies this note.
This work does not settle the full width-three or unrestricted conjecture.

\paragraph{Research and verification.}
Aafi Khan Malek directed the project. ChatGPT assisted with
proof development, drafting, formalization, and independent checking.
Codex was not used in this result.

\begin{table}[htbp]
\centering\footnotesize\setlength{\tabcolsep}{4pt}
\begin{tabular}{r l r @{\hspace{1em}} r l r}
\toprule
Row & Coefficients & $23B+20T$ & Row & Coefficients & $23B+20T$\\ \midrule
1 & $2\,\mathit{Xa}+3\,\mathit{Za}$ & 100 & 14 & $2\,\mathit{Ya}+\,\mathit{xa}$ & 60\\
2 & $\,\mathit{Xa}+\,\mathit{Ya}$ & 40 & 15 & $2\,\mathit{aX}+6\,\mathit{ay}+15\,\mathit{xa}$ & 460\\
3 & $30\,\mathit{ax}+60\,\mathit{ba}+5\,\mathit{xY}+3\,\mathit{xZ}+10\,\mathit{xy}$ & 2340 & 16 & $\,\mathit{Xa}$ & 20\\
4 & $4\,\mathit{ay}+6\,\mathit{xa}+\,\mathit{yX}$ & 220 & 17 & $2\,\mathit{Xa}+\,\mathit{ya}$ & 60\\
5 & $\,\mathit{xy}+7\,\mathit{ya}$ & 160 & 18 & $7\,\mathit{Xa}+\,\mathit{yX}$ & 160\\
6 & $2\,\mathit{ay}+4\,\mathit{ba}+\,\mathit{yx}$ & 152 & 19 & $7\,\mathit{Za}+\,\mathit{xy}$ & 160\\
7 & $20\,\mathit{ax}+5\,\mathit{xY}+3\,\mathit{xZ}+30\,\mathit{ya}$ & 1160 & 20 & $7\,\mathit{Xa}+\,\mathit{yx}$ & 160\\
8 & $7\,\mathit{xa}+\,\mathit{yx}$ & 160 & 21 & $7\,\mathit{Za}+\,\mathit{xY}$ & 160\\
9 & $6\,\mathit{ay}+12\,\mathit{ba}+\,\mathit{yX}+2\,\mathit{yx}$ & 456 & 22 & $3\,\mathit{YX}+5\,\mathit{Yx}+20\,\mathit{aY}+30\,\mathit{ya}$ & 1160\\
10 & $\,\mathit{Ya}$ & 20 & 23 & $7\,\mathit{Za}+\,\mathit{xZ}$ & 160\\
11 & $7\,\mathit{Ya}+\,\mathit{xy}$ & 160 & 24 & $\,\mathit{Za}$ & 20\\
12 & $7\,\mathit{Ya}+\,\mathit{xY}$ & 160 & 25 & $3\,\mathit{Za}+\,\mathit{xa}$ & 80\\
13 & $4\,\mathit{Yx}+10\,\mathit{aY}+15\,\mathit{ya}$ & 580 &  &  & \\
\bottomrule
\end{tabular}
\caption{Integer certificate rows. The independent audit enumerates the 70 components and 60 orders, checks all row inequalities and budgets, and checks coverage.}
\label{tab:cert}
\end{table}


\begin{thebibliography}{9}
\bibitem{sah} Ashwin Sah.
Improving the $1/3$--$2/3$ conjecture for width two posets.
\emph{Combinatorica} 41 (2021), 99--126.
\href{https://doi.org/10.1007/s00493-020-4091-3}{doi:10.1007/s00493-020-4091-3}.
\end{thebibliography}
\end{document}
